% 01_title.tex — Portada y declaraciones
% ════════════════════════════════════════════════════════════════════════════
% MEJORAS v2.0 (Opción B):
%   ✅ Título neutralizado (eliminada la frase que anticipaba hallazgo)
%   ✅ Título reducido a 21 palabras (ES) / 18 palabras (EN)
%   ✅ Co-autor añadido (placeholder [CO-REVISOR]) conforme a doble revisión
%   ✅ Placeholders prominentes para OSF y GitHub
% ════════════════════════════════════════════════════════════════════════════
\tituloprincipal{Aprendizaje automático y control avanzado para la regulación
del pH en la industria azucarera: una revisión de alcance y análisis
bibliométrico}{%
\textit{Machine Learning and Advanced Control for pH Regulation in the Sugar
Industry: A Scoping Review and Bibliometric Analysis}}

\vspace{0.6em}

\begin{center}
{\small \textbf{Abraham Castillo}$^{a,*}$~\orcidlink{0000-0001-6033-617X}
\quad\textbf{D. Barreto}$^{a}$~\orcidlink{0000-0002-0134-1291} \\[0.4em]
$^{a}$Escuela de Posgrado, Universidad Nacional de Ingeniería,
Av.\ Túpac Amaru 210, Lima 15333, Perú \\[0.3em]
$^{*}$\textit{Autor de correspondencia:}
\href{mailto:abraham.castillo.a@uni.pe}{abraham.castillo.a@uni.pe} \\[0.5em]
\textit{Recibido:} \hspace{2em} \textit{Revisado:} \hspace{2em}
\textit{Aceptado:} \hspace{2em} \textit{Publicado en línea:}}
\end{center}

\vspace{0.5em}

\begin{small}

\noindent\textbf{Declaración de financiamiento.}
Esta investigación no recibió financiamiento específico de agencias públicas,
del sector comercial ni de organismos sin fines de lucro.

\medskip
\noindent\textbf{Declaración de conflictos de interés.}
Los autores declaran que no existe ningún conflicto de interés conocido que
haya podido influir en los resultados reportados en este artículo.

\medskip
\noindent\textbf{Disponibilidad de datos y código.}
El protocolo, las cadenas booleanas por base (incluyendo SciELO y LILACS), los
registros de decisiones de cribado (humanas y automáticas), la matriz de
co-ocurrencia de palabras clave, la matriz de confusión inter-humano, los
resultados de los análisis de sensibilidad y el código del \textit{pipeline}
(cribado, verificación léxica y chequeo con modelo de lenguaje, con el
\textit{prompt} y el resumen SHA-256 del modelo) están disponibles en el Open
Science Framework (\href{\URLOSF}{\OSFTexto}) y en GitHub
(\href{\URLGitHub}{\GitHubTexto}).

\medskip
\noindent\textbf{Registro del protocolo.}
El protocolo de la revisión de alcance se registró en OSF de forma
\textbf{retrospectiva}, el \FechaRegistroOSF, es decir, después de ejecutar la
búsqueda (\FechaBusqueda) y de un cribado preliminar; por tanto, el registro
documenta el diseño pero no constituye una preespecificación independiente. Se
elaboró siguiendo la extensión PRISMA-ScR \parencite{tricco2018prismascr} y el
marco metodológico para revisiones de alcance
\parencite{arksey2005scoping,munn2018scoping}:
\href{\URLOSF}{\OSFTexto}. Las desviaciones respecto del diseño inicial se
declaran en la sección~\ref{sub:diseno}.

\medskip
\noindent\textbf{Aprobación ética.}
Esta investigación es una revisión de literatura publicada y no involucra
participantes humanos ni animales; por tanto, no requiere aprobación de
comité de ética.

\end{small}

\vspace{0.5em}
\hrule
\vspace{0.5em}

% ════════════════════════════════════════════════════════════════════════════
% TODO crítico antes del envío:
%   1. Sustituir [CO-REVISOR] por el nombre real del co-revisor identificado.
%   2. Sustituir [ORCID_COR] por el ORCID real (formato: 0000-0002-XXXX-XXXX).
%   3. Sustituir osf.io/[PENDIENTE] por la URL definitiva del registro OSF.
%   4. Sustituir github.com/[PENDIENTE] por la URL definitiva del repo público.
% ════════════════════════════════════════════════════════════════════════════
